% Source: ACM_Templates.pdf, PDF pp. 71-74; recursive Tarjan traversal.
% Review and validation: docs/validation.md.
\subsection{Tarjan 强连通分量与缩点}
\label{algo:scc}
\AlgorithmInfo
    {划分有向图中的强连通分量，并构造缩点后的 DAG。}
    {划分、缩点均为 $O(n+m)$ 时间，空间 $O(n+m)$。}
    {顶点与分量从 1 开始；支持自环、重边和不连通图。}

\subsubsection*{原理}
同一强连通分量内的点两两可达。Tarjan 维护发现时间 $dfn_u$、能回到的最早栈内位置 $low_u$，以及尚未归属某个分量的点栈。

搜索树边 $u\to v$ 完成后，用 $low_v$ 更新 $low_u$；遇到已发现且仍在点栈内的 $v$，用 $dfn_v$ 更新。已经分配到某个 SCC 的点不能再用来降低当前 $low$。若 $low_u=dfn_u$，从点栈弹到 $u$ 为止，这些点恰好形成一个 SCC。

模板使用递归 DFS，最坏深度 $O(n)$；\InlineCode{stk} 只保存尚未归属分量的顶点。

\subsubsection*{缩点与 DAG 上的计算}
将分量内部的边忽略，其余边按分量编号连接，得到有向无环图；若缩点图仍有环，环上各点本应属于同一个 SCC，矛盾。模板按弹栈顺序编号，因此跨分量边满足
\[
    bel_u>bel_v\qquad(u\to v,\ bel_u\ne bel_v).
\]
从大到小枚举分量编号就是一种拓扑序。\InlineCode{dag()} 保留平行边以维持线性复杂度；求入度时要按相同边集计数，若题目需要按不同邻点计数，再自行去重。

一般 DAG 的拓扑排序、路径计数和最短 / 最长路模板见第~\ref{algo:dag-dp}~节。

若题意允许收集分量内全部点权，可先求分量权和，再按第~\ref{algo:dag-dp}~节做点权 DP；起点限制和负权初始化沿用该节约定。

\subsubsection*{使用教程}
\InlineCode{init(n)} 后用 \InlineCode{add(u,v)} 加有向边。\InlineCode{build()} 返回分量数，\InlineCode{bel[u]} 是 $1\ldots cnt$ 内的分量编号，\InlineCode{sz[c]} 是分量大小；\InlineCode{dag()} 返回以分量编号为点的邻接表，保留平行边。新增边后重新构建即可，原图保留。
\begin{lstlisting}[style=mtycpp,numbers=none]
SCC g;
g.init(4);
g.add(1,2),g.add(2,1),g.add(2,3),g.add(3,4);
int cnt=g.build();
bool same=g.bel[1]==g.bel[2];
auto dag=g.dag();
\end{lstlisting}

\needspace{7\baselineskip}
\subsubsection*{模板代码}
\CodeFile{code/graphs/scc.cpp}
